%=======================================================================
%  CHAPTER 2 — NUMERICAL PROBLEMS AND SOLUTIONS
%=======================================================================
\section*{\color{acc}Ch 2 \textperiodcentered{} Cellular Concepts \gap
\textcolor{sub}{Numerical Problems \& Solutions}}
\vspace{-4pt}{\color{acc}\rule{\textwidth}{1.1pt}}\vspace{4pt}

{\color{sub}\footnotesize Every problem below is a real PYQ, with its year tags and marks.
Solutions are worked in full: what was calculated, why, and which formula was used.
\mk{Nothing is skipped to save space.}}

\lead{The five formulas everything here is built on}
\begin{itemize}
  \item \textbf{Cluster size} \; $N = i^2 + ij + j^2$ \; (allowed: 1, 3, 4, 7, 9, 12, 13, 16, 19, 21)
  \item \textbf{Co-channel reuse ratio} \; $Q = D/R = \sqrt{3N}$
  \item \textbf{Worst-case SIR} \; $\dfrac{S}{I} = \dfrac{(\sqrt{3N})^{\,n}}{i_0}$, \;
        so \; $N = \dfrac{1}{3}\left(i_0 \cdot \dfrac{S}{I}\right)^{2/n}$
  \item \textbf{Hexagonal cell area} \; $A = \dfrac{3\sqrt{3}}{2}R^2 = 2.598\,R^2$
  \item \textbf{Traffic} \; $A_u = \lambda H$, \; $A = U A_u$, \; and Erlang B for the GoS
\end{itemize}
{\color{sub}\footnotesize $i_0$ $=$ number of first-tier interferers:
\textbf{6} omni, \textbf{2} for $120^\circ$ sectoring, \textbf{1} for $60^\circ$ sectoring.
$n$ $=$ path loss exponent.}

\hr

\T{Problem 1 \quad Prove $Q=\sqrt{3N}$ and find the optimal $N$}
{\color{sub}\footnotesize \tS{6} \yr{\textbf{71 Bh} \m{4+4}, \texttt{79 Bh} \m{3+5},
\textbf{\texttt{80 Bh}} \m{4}, \textbf{70 Bh} \m{3+4+3}, \textbf{\texttt{81 Bh}} \m{3+3},
74 Ma \m{4}}}

\lead{1.1 \; SIR $=$ 15 dB worst case, $n = 4$, considering trunking efficiency
\yr{(\textbf{71 Bh})}}
\begin{asked}{2071 Bhadra \textperiodcentered{} Q2 \textperiodcentered{} [4+4]}
Prove that for a hexagonal geometry the co-channel reuse ratio is given by
$Q = \sqrt{3N}$; Where $N = i^2 + j^2 + ij$. A cellular service provider decides to use a
digital TDMA scheme which can tolerate a Signal-to-Interference Ratio of 15 dB in the worst
case. Find the optimal value of $N$ for
\begin{enumerate}[label=\alph*), leftmargin=1.6em]
  \item Omni directional antennas
  \item $120^\circ$ Sectoring
  \item $60^\circ$ Sectoring
\end{enumerate}
\par\vspace{1pt}
[Use path loss exponent of 4 and consider trunking efficiency]
\end{asked}

\begin{itemize}
  \item \textbf{Step 1: convert SIR to a linear ratio.} The formula is in power ratio, not dB.
        \[ \frac{S}{I} = 10^{15/10} = 10^{1.5} = 31.6228 \]
  \item \textbf{Step 2: apply} $\;N = \frac{1}{3}\left(i_0 \cdot S/I\right)^{2/n}\;$
        for each antenna configuration. Equivalently find $Q$ first, then $N = Q^2/3$.
\end{itemize}

\par\vspace{2pt}
\begin{tabularx}{\textwidth}{@{}L{2.6cm}L{1.2cm}XXL{2.4cm}@{}}
\toprule
\hrow \thead{Configuration} & \thead{$i_0$} & \thead{$Q=(i_0\,S/I)^{1/n}$} &
\thead{$N = Q^2/3$} & \thead{Round up to allowed $N$} \\
\midrule
Omni-directional & 6 & $(6\times31.6228)^{1/4} = (189.74)^{0.25} = 3.7114$ & $3.7114^2/3 = 4.5915$
  & \mk{$N=7$} \\
$120^\circ$ sectoring & 2 & $(2\times31.6228)^{1/4} = (63.246)^{0.25} = 2.8201$ &
  $2.8201^2/3 = 2.6509$ & \mk{$N=3$} \\
$60^\circ$ sectoring & 1 & $(31.6228)^{1/4} = 2.3714$ & $2.3714^2/3 = 1.8745$ & \mk{$N=3$} \\
\bottomrule
\end{tabularx}

\begin{itemize}
  \item \textbf{Step 3: why round \emph{up}, not to the nearest.}
        $N$ must be one of $1,3,4,7,9,12,\dots$, and a \textbf{smaller} $N$ gives a
        \textbf{worse} SIR. Rounding down would break the 15 dB requirement.
        \mk{So always take the next allowed value above the computed one.}
  \begin{itemize}
    \item omni: $4.5915 \rightarrow$ next allowed is 7 (4 is allowed but $4 < 4.5915$),
    \item $120^\circ$: $2.6509 \rightarrow 3$,
    \item $60^\circ$: $1.8745 \rightarrow 3$.
  \end{itemize}
  \item \textbf{Step 4: comment on the results \mk{(the marks for ``considering trunking
        efficiency'')}.}
  \begin{itemize}
    \item Sectoring cuts the required cluster size from \textbf{7 to 3}, so each cell keeps
          \mk{$7/3 = 2.33$ times as many channels}.
    \item But $60^\circ$ sectoring gives the \textbf{same} $N=3$ as $120^\circ$ here, while
          splitting each cell's channels into \textbf{6} pools instead of 3.
    \item A pool of 6 smaller groups carries \textbf{less} traffic at the same GoS than 3 larger
          ones. \mk{So $60^\circ$ buys no extra capacity and loses more trunking efficiency.}
    \item \textbf{Conclusion: $120^\circ$ sectoring with $N=3$ is the best choice.}
  \end{itemize}
\end{itemize}

\lead{1.2 \; SIR $=$ 15 dB worst case, $n = 3$ \yr{(\texttt{79 Bh})}}
\begin{asked}{2079 Bhadra \textperiodcentered{} Q2 \textperiodcentered{} [3+5]}
Define Grade of Service (GoS) and explain how can it be measured in a `blocked call cleared'
type of trunking system. A cellular service provider decides to use a digital TDMA scheme that
can tolerate a signal-to-interference ratio of 15 dB in the worst case. The mobile radio
channel provided a propagation path loss exponent of $n = 3$. Find the optimal value of $N$ for
(a)~Omni-directional antennas, (b)~1200 sectoring, and Ic)~$60^\circ$ sectoring. Comment on
your results.
\par\vspace{1pt}
{\color{sub}``1200 sectoring'' and ``Ic)'' are the paper's own typos for $120^\circ$ and (c).}
\end{asked}

\begin{itemize}
  \item Same $S/I = 31.6228$, but now the exponent is $2/n = 2/3$.
\end{itemize}
\par\vspace{2pt}
\begin{tabularx}{\textwidth}{@{}L{2.6cm}L{1.2cm}XXL{2.4cm}@{}}
\toprule
\hrow \thead{Configuration} & \thead{$i_0$} & \thead{$Q=(i_0\,S/I)^{1/3}$} &
\thead{$N = Q^2/3$} & \thead{Round up} \\
\midrule
Omni & 6 & $(189.74)^{1/3} = 5.7462$ & $5.7462^2/3 = 11.0064$ & \mk{$N=12$} \\
$120^\circ$ & 2 & $(63.246)^{1/3} = 3.9842$ & $3.9842^2/3 = 5.2913$ & \mk{$N=7$} \\
$60^\circ$ & 1 & $(31.6228)^{1/3} = 3.1623$ & $3.1623^2/3 = 3.3333$ & \mk{$N=4$} \\
\bottomrule
\end{tabularx}
\begin{itemize}
  \item \textbf{Comment.} A \textbf{smaller path loss exponent means interference decays more
        slowly}, so co-channel cells must sit further apart and $N$ rises sharply
        ($7 \rightarrow 12$ for omni).
  \item \mk{Sectoring is therefore even more valuable in a low-$n$ environment}: it takes $N$ from
        12 down to 4, a $3\times$ capacity gain.
\end{itemize}

\lead{1.3 \; The AMPS case, SIR $=$ 18 dB, $n = 4$ \yr{(standard textbook check)}}
\begin{asked}[PROBLEM]{not a PYQ \textperiodcentered{} the standard AMPS check \added}
An AMPS operator requires a worst-case signal-to-interference ratio of 18 dB with a path loss
exponent $n = 4$. Find the smallest allowable cluster size $N$ for omni-directional antennas.
\end{asked}
\begin{itemize}
  \item $S/I = 10^{1.8} = 63.0957$.
  \item Omni: $Q = (6 \times 63.0957)^{1/4} = (378.57)^{0.25} = 4.4110$, so
        $N = 4.4110^2/3 = 6.4857 \rightarrow \boxed{N = 7}$.
  \item \mk{This is exactly why every early AMPS system used a 7-cell cluster.}
\end{itemize}

\lead{1.4 \; Exact Vs approximate SIR for $N=7$, $n=4$}
\begin{asked}[PROBLEM]{not a PYQ \textperiodcentered{} the follow-up every examiner likes \added}
For a 7-cell cluster with $n = 4$, compute the worst-case $S/I$ twice: first with the
approximation that all six first-tier interferers lie at the same distance $D$, and then
exactly, with a mobile at the cell edge seeing interferers at $D-R$, $D$, $D$, $D$, $D$ and
$D+R$. Compare the two, in dB.
\end{asked}
\begin{itemize}
  \item \textbf{Approximate} (all six interferers taken at distance $D$):
        \[ \frac{S}{I} = \frac{(\sqrt{3\times7})^{4}}{6} = \frac{4.5826^4}{6}
           = \frac{441.0}{6} = 73.50 \;=\; 18.66\ \text{dB} \]
  \item \textbf{Exact} (mobile at the cell edge, true distances
        $D-R,\,D,\,D,\,D,\,D,\,D+R$ with $Q = D/R = 4.5826$):
        \[ \frac{S}{I} = \frac{1}{2(Q-1)^{-4} + 2(Q+1)^{-4} + 2Q^{-4}}
           = 53.376 \;=\; 17.27\ \text{dB} \]
  \item \mk{The approximation is optimistic by about 1.4 dB.} Both fall short of the 18 dB AMPS
        target, which is the standard justification for sectoring.
\end{itemize}

\hr

\T{Problem 2 \quad Spectrum, duplex channels and channels per cell}
{\color{sub}\footnotesize \tP{1} \yr{\textbf{70 Bh} \m{3+4+3}}}

\begin{asked}{2070 Bhadra \textperiodcentered{} Q2 \textperiodcentered{} [3+4+3]}
Explain the difference between co-channel and adjacent channel interference. Prove that the
co-channel reuse ratio is given by $Q = \sqrt{3N}$, where $N = i^2 + ij + j^2$ is the cluster
size. If 20 MHz of total spectrum is allocated for a duplex (i.e.\ bidirectional) wireless
cellular system and each simplex (i.e.\ one-way) channel has 25 KHz of bandwidth, find
\begin{enumerate}[label=\alph*), leftmargin=1.6em]
  \item The number of duplex channels, and
  \item The total number of channels per cell, assuming a cluster size of $N = 4$.
\end{enumerate}
\end{asked}
{\color{sub}\footnotesize The $Q=\sqrt{3N}$ proof is Problem 1 above; only the spectrum
arithmetic is new here.}

\begin{itemize}
  \item \textbf{Step 1: what one duplex channel actually costs.}
        A duplex (two-way) channel needs \textbf{two} simplex channels, one for each direction.
        \[ B_{\text{duplex}} = 2 \times 25\ \text{kHz} = 50\ \text{kHz} \]
        \mk{Missing this factor of 2 is the single most common error in this question.}
  \item \textbf{Step 2: total duplex channels.}
        \[ N_{\text{total}} = \frac{20 \times 10^{6}}{50 \times 10^{3}}
           = \boxed{400\ \text{duplex channels}} \]
  \item \textbf{Step 3: channels per cell for $N = 4$.}
        \[ k = \frac{S}{N} = \frac{400}{4} = \boxed{100\ \text{channels per cell}} \]
  \item \textbf{Step 4: the other cluster sizes, for comparison.}
  \begin{itemize}
    \item $N = 7$: $\;400/7 = 57.14 \rightarrow 57$ channels per cell,
    \item $N = 12$: $\;400/12 = 33.33 \rightarrow 33$ channels per cell.
  \end{itemize}
  \item \mk{Always round \emph{down} here}: a fractional channel cannot be allocated.
\end{itemize}

\hr

\T{Problem 3 \quad Minimum co-channel distance from the cell area}
{\color{sub}\footnotesize \tP{1} \yr{\textbf{74 Bh} \m{4}}}

\begin{asked}{2074 Bhadra \textperiodcentered{} Q2 (b) \textperiodcentered{} [4]}
For a seven cell reuse pattern, find the minimum distance between centers of co-channel cells.
Area of each cell is uniform and is equal to 23 square km.
\end{asked}

\begin{itemize}
  \item \textbf{Step 1: get $R$ from the area.} The cell is a regular hexagon of
        ``radius'' $R$ (centre to vertex), so
        \[ A = \frac{3\sqrt{3}}{2}R^2 = 2.598\,R^2 \]
        \[ R = \sqrt{\frac{A}{2.598}} = \sqrt{\frac{23}{2.598}} = \sqrt{8.8530}
             = 2.9754\ \text{km} \]
  \item \textbf{Step 2: apply the reuse ratio.} For $N = 7$,
        \[ Q = \sqrt{3N} = \sqrt{21} = 4.5826 \]
  \item \textbf{Step 3: minimum co-channel distance.}
        \[ D = Q \cdot R = 4.5826 \times 2.9754 = \boxed{13.635\ \text{km}} \]
  \item \textbf{Sanity check.} $D/R = 13.635/2.9754 = 4.58$ \checkmark matches $\sqrt{21}$.
  \item \mk{If the question had given the area as a circle instead}, use $A = \pi R^2$ and
        $R = \sqrt{23/\pi} = 2.7052$ km, giving $D = 12.397$ km. Read the wording carefully:
        ``uniform cell'' in this paper means the hexagon.
\end{itemize}

\hr

\T{Problem 4 \quad Users supported and market penetration}
{\color{sub}\footnotesize \tF{3} \yr{\texttt{80 Ba}, \texttt{81 Ba} \m{5}, \textbf{79 Ch}}}

\begin{asked}{2080 Baishakh \textperiodcentered{} Q2 (a) \textperiodcentered{} [5]
\ \textperiodcentered{} \ 2079 Chaitra \textperiodcentered{} Q2 (a) \textperiodcentered{} [5]}
An urban area has a population of two million residents. System A has 394 cells with 19
channels each. Find the number of users that can be supported at 2\% blocking if each user
average two calls per hour at average call duration of three minutes. Assuming that trunk
systems are operated at maximum capacity, compute the percentage market penetration of
cellular provider.
\par\vspace{1pt}
{\color{sub}Both papers print this word for word (79 Ch reads ``averages two calls''), and
79 Ch attaches the Erlang B table with the $N = 19$ row: 11.23 at 1\%, \textbf{12.33 at 2\%}.}
\end{asked}

\begin{itemize}
  \item \textbf{Step 1: traffic intensity per user.} $\;A_u = \lambda H$, and both terms must be in
        the \emph{same} time unit.
        \[ \lambda = 2\ \text{calls/hour}, \qquad H = 3\ \text{min} = \frac{3}{60}\ \text{hour} \]
        \[ A_u = 2 \times \frac{3}{60} = \boxed{0.1\ \text{Erlang per user}} \]
  \item \textbf{Step 2: carried traffic per cell, from the Erlang B table.}
        With $C = 19$ channels and GoS $=$ 0.02,
        \[ A = 12\ \text{Erlang} \]
        {\color{sub}\footnotesize (The exact Erlang B value is 12.333 Erlang; the standard table
        rounds to 12, which is what the marking scheme expects.)}
  \item \textbf{Step 3: users per cell.}
        \[ U = \frac{A}{A_u} = \frac{12}{0.1} = 120\ \text{users per cell} \]
  \item \textbf{Step 4: total subscribers.}
        \[ U_{\text{total}} = 120 \times 394 = \boxed{47\,280\ \text{subscribers}} \]
  \item \textbf{Step 5: market penetration.}
        \[ \frac{47\,280}{2\,000\,000} = 0.02364 = \boxed{2.364\ \%} \]
  \item \mk{Note the assumption stated in the question}: ``trunked systems operate at maximum
        capacity''. That is what allows the full 12 Erlang per cell to be used.
  \item {\color{sub}\footnotesize Using the exact 12.333 Erlang instead gives 123.33 users/cell,
        48\,592 subscribers and 2.430\,\% penetration. Either is acceptable if the table value
        used is stated.}
\end{itemize}

\hr

\T{Problem 5 \quad Trunking with 1 and 5 channels, then doubled load}
{\color{sub}\footnotesize \tP{1} \yr{\textbf{77 Ch} \m{2+2+2+3}}}

\begin{asked}{2077 Chaitra \textperiodcentered{} Q2 \textperiodcentered{} [2+2+2+3]}
A mobile radio system where each user averages three calls per hour and each call lasting an
average of 5 minutes.
\begin{enumerate}[label=\alph*), leftmargin=1.6em]
  \item What is traffic intensity for each user?
  \item Find the number of users that could use the system with 1\% blocking if only one
        channel is available.
  \item Find the number of users that could use the system with 1\% blocking if five trunked
        channels are available.
  \item If the number of users you found in (c) is suddenly doubled, what is the new blocking
        probability of the five channel trunked mobile radio system? Justify whether the
        performance is acceptable or not.
\end{enumerate}
\par\vspace{2pt}
\begin{tabularx}{0.94\textwidth}{@{}L{2.9cm}XXXXXXXX@{}}
\toprule
\hrow \thead{No.\ of Trunks (N)} & \thead{0.1\%} & \thead{0.2\%} & \thead{0.5\%} & \thead{1\%}
 & \thead{1.2\%} & \thead{3\%} & \thead{5\%} & \thead{7\%} \\
1 & 0.001 & 0.002 & 0.005 & 0.010 & 0.012 & 0.031 & 0.053 & 0.075 \\
2 & 0.046 & 0.065 & 0.105 & 0.153 & 0.168 & 0.282 & 0.381 & 0.470 \\
3 & 0.194 & 0.249 & 0.349 & 0.455 & 0.489 & 0.715 & 0.899 & 1.06 \\
4 & 0.439 & 0.535 & 0.701 & 0.869 & 0.922 & 1.26 & 1.52 & 1.75 \\
5 & 0.762 & 0.900 & 1.13 & 1.36 & 1.43 & 1.88 & 2.22 & 2.50 \\
\bottomrule
\end{tabularx}
\par\vspace{1pt}
{\color{sub}Table caption as printed: ``Traffic (A) in erlangs for B\%''.}
\end{asked}

\begin{itemize}
  \item \textbf{(a) Traffic intensity per user.}
        \[ A_u = \lambda H = 3 \times \frac{5}{60} = \boxed{0.25\ \text{Erlang}} \]
  \item \textbf{(b) One channel, GoS $=$ 1\,\%.}
        From the printed table, $C = 1$ and $B = 1\,\%$ give $A = 0.010$ Erlang.
        \[ U = \frac{A}{A_u} = \frac{0.010}{0.25} = 0.04\ \text{users} \]
        \[ \Rightarrow \boxed{0\ \text{users}} \]
        \mk{Interpretation:} a single channel cannot support even one such user at 1\,\% blocking.
        A lone user offering 0.25 Erlang to one channel would be blocked about
        $0.25/(1+0.25) = 20\,\%$ of the time.
  \item \textbf{(c) Five trunked channels, GoS $=$ 1\,\%.}
        From the table, $C = 5$ and $B = 1\,\%$ give $A = 1.36$ Erlang.
        \[ U = \frac{1.36}{0.25} = 5.44 \;\Rightarrow\; \boxed{5\ \text{users}} \]
  \item \mk{This is the whole point of trunking:} 1 channel supports 0 users, but 5 channels
        support 5. \textbf{Capacity grows faster than the channel count.}
  \item \textbf{(d) Users doubled to 10.}
        \[ A = U A_u = 10 \times 0.25 = 2.5\ \text{Erlang} \]
        Apply the Erlang B formula with $C = 5$, $A = 2.5$:
        \[ \Pr[\text{blocking}] = \frac{A^C/C!}{\sum_{k=0}^{C} A^k/k!} \]
        \begin{center}\small
        \begin{tabular}{@{}lllllll@{}}
        \toprule
        $k$ & 0 & 1 & 2 & 3 & 4 & 5 \\
        \midrule
        $A^k/k!$ & 1 & 2.5 & 3.125 & 2.6042 & 1.6276 & 0.8138 \\
        \bottomrule
        \end{tabular}
        \end{center}
        \[ \sum_{k=0}^{5} \frac{A^k}{k!} = 11.6706, \qquad
           \Pr[\text{blocking}] = \frac{0.8138}{11.6706} = 0.0697 \]
        \[ \boxed{\text{GoS} = 6.97\ \% } \]
  \item \textbf{Justification.} The design target was 1\,\%. Doubling the users has pushed
        blocking to \mk{almost 7\,\%, seven times worse}.
        \textbf{This is not acceptable}: roughly 1 call in 14 fails in the busy hour.
  \item \mk{Note how non-linear this is}: doubling the load multiplied the blocking by 7, not by 2.
        That non-linearity is the reason cellular systems are dimensioned on the busy hour.
\end{itemize}

\hr

\T{Problem 6 \quad Full system design from area and spectrum}
{\color{sub}\footnotesize \tP{1} \yr{72 Ma \m{12}} \gap$\cdot$\gap \mk{the longest numerical in
the chapter}}

\begin{asked}{2072 Magh \textperiodcentered{} Q2 \textperiodcentered{} [12]}
A city with a coverage area of 1500 sq km is covered with a 12-cell system each with a radius
of 1.387 km. If the total spectrum allocated is 28.5 MHz with a full duplex channel bandwidth
of 25 MHz. Assume a GOS of 0.02 for a blocked calls cleared system, is specified and the
offered traffic per user is 0.03 Erlangs and traffic intensity of each cell is 84 Erlang,
compute:
\begin{enumerate}[label=(\alph*), leftmargin=1.9em]
  \item the number of cells in the service area
  \item the number of channels per cell
  \item the maximum carrier traffic
  \item the total number of users that can be served for 2\% GOS
  \item the number of mobiles per unique channel.
  \item Theoretical maximum number of users that could be served at one time by the system.
\end{enumerate}
\par\vspace{1pt}
{\color{sub}\mk{The paper prints ``25 MHz'' for the channel bandwidth}: a typo for
\textbf{25 kHz}, since 28.5 MHz would otherwise buy one channel. Solve with 25 kHz and say so.}
\end{asked}

\begin{itemize}
  \item \textbf{(a) Number of cells in the service area.}
        \[ A_{\text{cell}} = 2.598\,R^2 = 2.598 \times (1.387)^2 = 2.598 \times 1.9238
           = 4.998\ \text{km}^2 \]
        \[ n_{\text{cells}} = \frac{1500}{4.998} = 300.12 \;\Rightarrow\;
           \boxed{300\ \text{cells}} \]
  \item \textbf{(b) Channels per cell.}
        \[ N_{\text{total}} = \frac{28.5 \times 10^{6}}{25 \times 10^{3}}
           = 1140\ \text{duplex channels} \]
        {\color{sub}\footnotesize The question says ``full duplex channel bandwidth 25 kHz'', so
        the $\times 2$ of Problem 2 does \emph{not} apply here. Read the wording.}
        \[ k = \frac{1140}{12} = \boxed{95\ \text{channels per cell}} \]
  \item \textbf{(c) Maximum carried traffic per cell.}
        The question supplies this as \textbf{84 Erlang}.
        {\color{sub}\footnotesize Cross-check from Erlang B: $C = 95$ at GoS $=$ 0.02 gives
        $A = 83.13$ Erlang, so the given 84 is consistent.}
        \[ A_{\text{cell}} = \boxed{84\ \text{Erlang}} \]
        System total $= 84 \times 300 = 25\,200$ Erlang.
  \item \textbf{(d) Total users served at 2\,\% GoS.}
        \[ U_{\text{cell}} = \frac{A_{\text{cell}}}{A_u} = \frac{84}{0.03} = 2800
           \ \text{users per cell} \]
        \[ U_{\text{total}} = 2800 \times 300 = \boxed{840\,000\ \text{users}} \]
  \item \textbf{(e) Mobiles per unique channel.}
        \[ \frac{U_{\text{cell}}}{k} = \frac{2800}{95} = \boxed{29.47\ \text{mobiles per channel}} \]
        \mk{Ie roughly 29 subscribers share each radio channel}, which is exactly the statistical
        multiplexing gain trunking delivers.
  \item \textbf{(f) Theoretical maximum users served at one time.}
        Every channel in every cell busy simultaneously:
        \[ 95 \times 300 = \boxed{28\,500\ \text{simultaneous calls}} \]
  \item \mk{Sanity check on (d) against (f):} 840\,000 subscribers share 28\,500 simultaneous
        circuits, a ratio of 29.5:1, matching (e). \checkmark
\end{itemize}

\hr

\T{Problem 7 \quad Area, channels and simultaneous calls}
{\color{sub}\footnotesize \tP{1} \yr{\textbf{\texttt{82 Bh}} \m{4}}}

\begin{asked}{2082 Bhadra \textperiodcentered{} Q2 (b) \textperiodcentered{} [4]}
The cellular system has 32 cells, each cell has a 1.6 km radius and the system reuse factor is
7. The system is to support 336 traffic channels in total. Determine the total geographical
area covered, the number of traffic channels per cell and the total number of simultaneous
calls supported by this system.
\end{asked}

\begin{itemize}
  \item \textbf{Step 1: area of one cell.}
        \[ A_{\text{cell}} = 2.598 \times (1.6)^2 = 2.598 \times 2.56 = 6.6509\ \text{km}^2 \]
  \item \textbf{Step 2: total geographical area.}
        \[ A_{\text{total}} = 32 \times 6.6509 = \boxed{212.83\ \text{km}^2} \]
  \item \textbf{Step 3: channels per cell.}
        \[ k = \frac{S}{N} = \frac{336}{7} = \boxed{48\ \text{channels per cell}} \]
  \item \textbf{Step 4: simultaneous calls.}
        Every cell can run all 48 of its channels at once, and there are 32 cells:
        \[ 48 \times 32 = \boxed{1536\ \text{simultaneous calls}} \]
  \item \mk{Watch the trap:} the 336 channels are the \textbf{system} total, reused by every
        cluster. The system supports 1536 calls, not 336, because the cluster repeats
        $32/7 = 4.57$ times.
\end{itemize}

\hr

\T{Problem 8 \quad Trunking loss from omni to $60^\circ$ sectoring}
{\color{sub}\footnotesize \tP{1} \yr{\texttt{81 Ba} \m{2+2+2}} \gap$\cdot$\gap
\mk{the one question where sectoring is shown to \emph{cost} capacity}}

\begin{asked}{2081 Baishakh \textperiodcentered{} Q2 (a) \textperiodcentered{} [2+2+2]}
For a cluster $N = 7$ system with a blocking probability of 1\% and an average call length of
two minutes, find the traffic capacity loss for a blocked calls cleared system due to trunking
for 57 channels when going from omnidirectional antennas to $60^\circ$ sectored antennas. The
average request rate per user is 1 call per hour.
\par\vspace{2pt}
\begin{tabularx}{0.62\textwidth}{@{}L{3.4cm}XXX@{}}
\toprule
\hrow \thead{GOS / Channel} & \thead{0.5\%} & \thead{1\%} & \thead{2\%} \\
1  & .005  & .0101 & .0204 \\
9  & 3.333 & 3.783 & 4.345 \\
19 & 10.33 & 11.23 & 12.33 \\
54 & 39.47 & 41.51 & 44.00 \\
57 & 42.1  & 44.2  & 46.8  \\
\bottomrule
\end{tabularx}
\end{asked}

\begin{itemize}
  \item \textbf{Step 1: traffic per user.}
        \[ A_u = \lambda H = 1 \times \frac{2}{60} = 0.0333\ \text{Erlang} \]
  \item \textbf{Step 2: omni-directional case.}
        All 57 channels sit in \textbf{one} trunked pool.
        From the table, $C = 57$ at 1\,\% gives
        \[ A_{\text{omni}} = 44.2\ \text{Erlang} \]
        \[ U_{\text{omni}} = \frac{44.2}{0.0333} = 1326\ \text{users per cell} \]
  \item \textbf{Step 3: $60^\circ$ sectored case.}
        The cell is split into \textbf{6 sectors}, and each sector gets its own channel group:
        \[ \frac{57}{6} = 9.5 \;\Rightarrow\; 9\ \text{channels per sector} \]
        \mk{Round down: a channel cannot be split between sectors.}
        From the table, $C = 9$ at 1\,\% gives $A = 3.783$ Erlang \textbf{per sector}, so
        \[ A_{\text{sectored}} = 6 \times 3.783 = 22.698\ \text{Erlang} \]
        \[ U_{\text{sectored}} = \frac{22.698}{0.0333} = 681\ \text{users per cell} \]
  \item \textbf{Step 4: the loss.}
        \[ \Delta A = 44.2 - 22.698 = \boxed{21.50\ \text{Erlang lost}} \]
        \[ \frac{21.50}{44.2} \times 100 = \boxed{48.65\ \% \text{ capacity loss}} \]
  \item \textbf{Why this happens \mk{(state this, it is where the marks are)}.}
  \begin{itemize}
    \item Erlang B is \textbf{strongly non-linear} in the channel count. One pool of 57 channels
          carries 44.2 Erlang, but six pools of 9 carry only $6 \times 3.783 = 22.7$ Erlang.
    \item That is the \textbf{loss of trunking efficiency}: a large pool absorbs random peaks far
          better than several small ones.
    \item Three channels are also lost outright to the rounding ($57 \rightarrow 6 \times 9 = 54$).
  \end{itemize}
  \item \mk{The complete engineering picture:} sectoring \textbf{improves SIR}, which lets $N$ be
        reduced, which gives \emph{more} channels per cell. That gain must beat the
        \textbf{48.65\,\% trunking loss} computed here, otherwise sectoring is a net loss. This
        problem deliberately shows only the cost side.
\end{itemize}

\hr

\T{Problem 9 \quad Calls per cell per hour, and SIR}
{\color{sub}\footnotesize \yr{lecture-note example, same pattern as \textbf{79 Bh} and 72 Ma}}

\begin{asked}[PROBLEM]{not a PYQ \textperiodcentered{} lecture-note example \added}
A cellular system has 416 radio channels available, of which 21 are control channels. The
average channel holding time is 3 minutes and the blocking probability during the busy hour is
2\%. The frequency reuse factor is 9. Find (i)~the number of calls per cell per hour, and
(ii)~the signal to co-channel interference ratio in dB (take $n = 4$).
\end{asked}

\begin{itemize}
  \item \textbf{Step 1: voice channels.}
        \[ 416 - 21 = 395\ \text{voice channels} \]
  \item \textbf{Step 2: voice channels per cell.}
        \[ \frac{395}{9} = 43.89 \;\Rightarrow\; 44\ \text{channels per cell}
           \ \text{(taken as 44 for the table lookup)} \]
  \item \textbf{Step 3: carried traffic.}
        From the Erlang B table, $C = 44$ at GoS $=$ 0.02 gives
        \[ A = 34.683\ \text{Erlang} \]
  \item \textbf{Step 4: convert Erlang to calls.}
        1 Erlang $=$ one continuously occupied channel. With a 3-minute holding time, one
        Erlang carries $1/3$ call per minute.
        \[ \text{calls per cell per minute} = \frac{34.683}{3} = 11.561 \]
        \[ \text{calls per cell per hour} = 11.561 \times 60 = \boxed{694\ \text{calls/cell/hour}} \]
  \item \textbf{Step 5: SIR for $N = 9$, $n = 4$.}
        \[ Q = \sqrt{3 \times 9} = \sqrt{27} = 5.1962 \]
        \[ \frac{S}{I} = \frac{Q^4}{6} = \frac{729}{6} = 121.5 \]
        \[ \frac{S}{I}\,[\text{dB}] = 10\log_{10}(121.5) = \boxed{20.85\ \text{dB}} \]
  \item \mk{Comment:} 20.85 dB comfortably exceeds the 18 dB AMPS requirement, so $N=9$ is a
        conservative, high-quality plan. $N=7$ would give 18.66 dB, barely adequate.
\end{itemize}

\hr

\T{Problem 10 \quad Cell splitting: the transmit power reduction}
{\color{sub}\footnotesize \tF{3} \yr{\textbf{73 Bh}, \textbf{75 Bh}, \textbf{\texttt{81 Bh}}}
\gap$\cdot$\gap usually a 1 or 2 mark part inside the sectoring Vs splitting question}

\begin{asked}{2081 Bhadra \textperiodcentered{} Q3 \textperiodcentered{} [4]
\ \textperiodcentered{} \ 2075 Bhadra \textperiodcentered{} Q3 \textperiodcentered{} [5]}
A telephone network company needs to expand its capacity based on demand on a city. A group of
engineers was selected to find the solution. Among the solution sectoring and cell splitting
were major technique for expansion purpose. Being a cellular planning engineer, which option do
you think is best and why?
\end{asked}
\begin{asked}{2073 Bhadra \textperiodcentered{} Q2 (a) \textperiodcentered{} [5]}
Explain how cell splitting and sectoring improve coverage and capacity in cellular system?
\end{asked}
{\color{sub}\footnotesize No paper prints numbers for this. The transmit-power result below is
the only computable part, and it is what turns a prose answer into a full-mark one. \added}

\begin{itemize}
  \item \textbf{Requirement.} The new microcell of radius $R/2$ must deliver the \textbf{same
        received power at its own boundary} as the old cell did at its boundary. Only then does
        the reuse plan carry over unchanged.
  \item \textbf{Step 1: write both boundary powers.}
        \[ P_r[\text{old boundary}] \propto P_{t1} R^{-n}, \qquad
           P_r[\text{new boundary}] \propto P_{t2} \left(\frac{R}{2}\right)^{-n} \]
  \item \textbf{Step 2: set them equal.}
        \[ P_{t1} R^{-n} = P_{t2} \left(\frac{R}{2}\right)^{-n}
           \;\Longrightarrow\; \frac{P_{t2}}{P_{t1}} = \left(\frac{1}{2}\right)^{n} \]
  \item \textbf{Step 3: put $n = 4$.}
        \[ \frac{P_{t2}}{P_{t1}} = \frac{1}{16} = 0.0625 \]
        \[ 10\log_{10}(0.0625) = \boxed{-12.04\ \text{dB}, \text{ ie reduce power by 12 dB}} \]
  \item \textbf{Capacity effect.} Halving $R$ needs about \textbf{4 times as many cells}
        ($A \propto R^2$), so the number of clusters, and therefore the capacity, rises about
        \mk{4 times} for the same spectrum.
\end{itemize}

\hr

\T{Erlang B reference table}
{\color{sub}\footnotesize Values recomputed from
$\Pr[\text{blocking}] = \dfrac{A^C/C!}{\sum_{k=0}^{C}A^k/k!}$ and cross-checked against the
tables printed on the 77 Ch and 81 Ba papers. \mk{Both printed tables agree with this one} once
their differing row labels (``Trunks $N$'' Vs ``GOS / Channel'') are read correctly.}

\penalty0\vspace{2pt}
\begin{center}\small
\begin{tabular}{@{}rrrrr@{}}
\toprule
\hrow \thead{Channels $C$} & \thead{$A$ at 0.5\,\%} & \thead{$A$ at 1\,\%} &
\thead{$A$ at 2\,\%} & \thead{$A$ at 5\,\%} \\
\midrule
1  &  0.005 &  0.010 &  0.020 &  0.053 \\
2  &  0.105 &  0.153 &  0.223 &  0.381 \\
3  &  0.349 &  0.455 &  0.602 &  0.899 \\
4  &  0.701 &  0.869 &  1.092 &  1.525 \\
5  &  1.132 &  1.361 &  1.657 &  2.219 \\
9  &  3.333 &  3.783 &  4.345 &  5.370 \\
10 &  3.961 &  4.461 &  5.084 &  6.216 \\
19 & 10.331 & 11.230 & 12.333 & 14.315 \\
20 & 11.092 & 12.031 & 13.182 & 15.249 \\
44 & 30.797 & 32.543 & 34.682 & 38.557 \\
54 & 39.474 & 41.505 & 43.997 & 48.536 \\
57 & 42.109 & 44.222 & 46.816 & 51.548 \\
95 & 76.325 & 79.368 & 83.133 & 90.123 \\
\bottomrule
\end{tabular}
\end{center}

\penalty0\vspace{2pt}
{\color{sub}\footnotesize \mk{How to use it in the exam:} you are always given either $C$ and the
GoS (look up $A$, then divide by $A_u$ for users), or $C$ and $A$ (look across the row for the
GoS). Never try to sum the series by hand unless $C \le 5$.}
